\section{Discussion}
\label{sec:discussion}

\paragraph{Information is usable relative to a reader.}
The BIOT host of Section~\ref{sec:rotation} is the clearest demonstration of the quantity the hypothesis is built on. The same isometric rotation that is neutral in the factorized reader adds 0.20 $\kappa$ on TUEV and removes 0.17 AUC-PR on CHB-MIT inside a reader that flattens bands into channels. For that reader a cross-band rotation is a feature it cannot undo, helpful where the label is envelope-dominated and harmful where morphology must be preserved. The consequence for the tokenizer literature is methodological: a transplant experiment that holds each host's reader fixed cannot by itself attribute a gain to content, and a tokenizer effect measured under one reader does not transfer to another. It also explains the one result that could be read as a counterexample to the hypothesis, CBraMod's stem scoring 0.738 under our reader: the stem exposes a time-frequency magnitude, which is the first clause of P8.

\paragraph{Explicit cross-frequency content is not required.}
Periodic discharges are fast activity locked to a slow rhythm, and it is natural to expect that a token which carries the measured locking should help. The archive says otherwise, and the hypothesis says why: the labels are functionals of envelope and waveform, the reader attends across bands itself, and an isometric rotation by the measured direction adds no usable information while a variant that makes the coupling readable alters the content magnitudes the labels depend on and loses 0.033 $\kappa$. Cole and Voytek's observation~\cite{cole2017brain} closes the loop: the non-sinusoidal spike-and-wave morphology that the waveform view carries is precisely the structure that inflates coupling estimates, so a model that profits from morphology will always look as if it profits from coupling unless the control is a matched untransformed token rather than a waveform-only one.

\paragraph{Why pretraining did not help.}
Amplitude-reconstruction pretraining helped on three corpora with scarce or noisy labels and hurt on nine with plentiful ones, and a frozen probe of the pretrained representation was the best probe on TUEV while fine-tuning from it was worse than from scratch. Both are expected when the label-relevant content is already explicit in the tokens: the reconstruction objective, dominated by $1/f$ amplitude~\cite{kommineni2026aperiodic,yu2026fame}, optimises something the task does not need, and the supervised recipe finds the usable information in one or two epochs without it.

\paragraph{Scope.}
The hypothesis predicts gains where labels are band-specific sub-second events and no gain where they are 30\,s spectral states or cohort-level diagnoses. The archive matches: TUEV, TUSZ, IIIC and TUEP lead; CAUEEG and ADFD trail; the single-seed Sleep-EDF and TUAR leads under the 16-band recipe are unexpected and must be confirmed before they are claimed. The three-axis advantage over two-axis (seven of eight corpora) is consistent with Section~\ref{sec:factorized} but confounded with a 33\% parameter difference until the matched control of Section~\ref{sec:limitations} is run.
